\documentclass[11pt,a4paper]{article}
\usepackage[utf8]{inputenc}
\usepackage[french]{babel}
\usepackage{amsmath}
\usepackage{amssymb}
\usepackage{geometry}
\usepackage{enumitem}
\usepackage{xcolor}
\usepackage{titlesec}
\usepackage{fancyhdr}
\usepackage{booktabs}
\usepackage{array}
\usepackage{listings}
\usepackage{tcolorbox}
\usepackage{algorithm}
\usepackage{algorithmic}

% Configuration de la page
\geometry{margin=2.5cm}
\pagestyle{fancy}
\fancyhf{}
\fancyhead[L]{\leftmark}
\fancyhead[R]{Solutions Machine Learning}
\fancyfoot[C]{\thepage}

% Configuration des titres
\titleformat{\section}
{\Large\bfseries\color{blue!70!black}}
{}
{0em}
{}[\titlerule]

\titleformat{\subsection}
{\large\bfseries\color{blue!50!black}}
{}
{0em}
{}

% Configuration pour le code Python
\lstset{
    language=Python,
    basicstyle=\ttfamily\small,
    keywordstyle=\color{blue}\bfseries,
    commentstyle=\color{gray}\itshape,
    stringstyle=\color{red},
    numbers=left,
    numberstyle=\tiny\color{gray},
    stepnumber=1,
    numbersep=5pt,
    backgroundcolor=\color{gray!10},
    frame=single,
    breaklines=true,
    showstringspaces=false,
    tabsize=4,
    morekeywords={np, pd, sklearn, plt, sns},
    inputencoding=utf8,
    extendedchars=true,
    literate={à}{{\`a}}1 {é}{{\'e}}1 {è}{{\`e}}1 {ê}{{\^e}}1 {ë}{{\"e}}1
             {ù}{{\`u}}1 {û}{{\^u}}1 {ü}{{\"u}}1 {ô}{{\^o}}1 {ç}{{\c c}}1
             {À}{{\`A}}1 {É}{{\'E}}1 {È}{{\`E}}1 {Ê}{{\^E}}1 {Ë}{{\"E}}1
             {Ù}{{\`U}}1 {Û}{{\^U}}1 {Ü}{{\"U}}1 {Ô}{{\^O}}1 {Ç}{{\c C}}1
}

% Configuration pour les zones importantes
\tcbuselibrary{breakable}
\newtcolorbox{solutionbox}{
    colback=green!5,
    colframe=green!50,
    breakable,
    title=Solution
}

\newtcolorbox{tipbox}{
    colback=yellow!5,
    colframe=yellow!50,
    breakable,
    title=Tip
}

\newtcolorbox{importantbox}{
    colback=red!5,
    colframe=red!50,
    breakable,
    title=Important
}

\newtcolorbox{formulebox}{
    colback=blue!5,
    colframe=blue!50,
    breakable,
    title=Formule
}

% Métadonnées
\title{Solutions - Exercices Machine Learning\\
\large Régression Linéaire, KNN, K-means, Arbre de Décision}
\author{AmineGR03}
\date{\today}

\begin{document}

\maketitle

\tableofcontents
\newpage

% ============================================
% PARTIE 1 : RÉGRESSION LINÉAIRE
% ============================================

\section{Régression Linéaire}

\subsection{Solution Exercice 1 : Calcul de la MSE et des paramètres}

\begin{solutionbox}
\textbf{Données :}
\begin{center}
\begin{tabular}{|c|c|c|c|c|}
\hline
\textbf{X} & \textbf{Y} & \textbf{XY} & \textbf{X²} & \textbf{Y²} \\
\hline
1 & 45000 & 45000 & 1 & 2025000000 \\
2 & 50000 & 100000 & 4 & 2500000000 \\
3 & 60000 & 180000 & 9 & 3600000000 \\
4 & 80000 & 320000 & 16 & 6400000000 \\
5 & 110000 & 550000 & 25 & 12100000000 \\
\hline
\textbf{$\sum$} & \textbf{345000} & \textbf{1195000} & \textbf{55} & \textbf{26625000000} \\
\hline
\end{tabular}
\end{center}

\textbf{1. Calcul de $a$ et $b$ :}

$n = 5$, $\sum X = 15$, $\sum Y = 345000$, $\sum XY = 1195000$, $\sum X^2 = 55$

\begin{align*}
a &= \frac{n\sum XY - \sum X \sum Y}{n\sum X^2 - (\sum X)^2} \\
  &= \frac{5 \times 1195000 - 15 \times 345000}{5 \times 55 - 15^2} \\
  &= \frac{5975000 - 5175000}{275 - 225} \\
  &= \frac{800000}{50} \\
  &= 16000
\end{align*}

\begin{align*}
b &= \frac{\sum Y - a\sum X}{n} \\
  &= \frac{345000 - 16000 \times 15}{5} \\
  &= \frac{345000 - 240000}{5} \\
  &= \frac{105000}{5} \\
  &= 21000
\end{align*}

\textbf{Modèle :} $\hat{Y} = 16000X + 21000$

\textbf{2. Calcul des prédictions :}

\begin{center}
\begin{tabular}{|c|c|c|c|}
\hline
\textbf{X} & \textbf{Y réel} & \textbf{$\hat{Y}$} & \textbf{Erreur (Y - $\hat{Y}$)} \\
\hline
1 & 45000 & 37000 & 8000 \\
2 & 50000 & 53000 & -3000 \\
3 & 60000 & 69000 & -9000 \\
4 & 80000 & 85000 & -5000 \\
5 & 110000 & 101000 & 9000 \\
\hline
\end{tabular}
\end{center}

\textbf{3. Calcul de la MSE :}

\begin{align*}
MSE &= \frac{1}{n}\sum_{i=1}^{n}(Y_i - \hat{Y}_i)^2 \\
    &= \frac{1}{5}[(8000)^2 + (-3000)^2 + (-9000)^2 + (-5000)^2 + (9000)^2] \\
    &= \frac{1}{5}[64000000 + 9000000 + 81000000 + 25000000 + 81000000] \\
    &= \frac{260000000}{5} \\
    &= 52000000
\end{align*}

\textbf{4. Interprétation :}
\begin{itemize}
    \item \textbf{$a = 16000$} : Chaque année d'expérience supplémentaire augmente le salaire de 16000€ en moyenne.
    \item \textbf{$b = 21000$} : Le salaire de base (sans expérience) est de 21000€.
\end{itemize}
\end{solutionbox}

\begin{tipbox}
\textbf{Tips pour le calcul :}
\begin{itemize}
    \item Vérifiez toujours que $\sum X^2 - (\sum X)^2/n \neq 0$ (sinon division par zéro)
    \item Organisez vos calculs dans un tableau pour éviter les erreurs
    \item La MSE est toujours positive (somme de carrés)
    \item Plus la MSE est petite, meilleur est le modèle
\end{itemize}
\end{tipbox}

\subsection{Solution Exercice 2 : Gradient Descent}

\begin{solutionbox}
\textbf{Données :} $a_0 = 0$, $b_0 = 0$, $\alpha = 0.01$, $n = 3$

Points : $(2, 4)$, $(3, 6)$, $(4, 8)$

\textbf{1. Calcul des prédictions initiales :}

Pour chaque point : $\hat{Y}_i = a_0 \times X_i + b_0 = 0 \times X_i + 0 = 0$

\begin{center}
\begin{tabular}{|c|c|c|c|}
\hline
\textbf{X} & \textbf{Y} & \textbf{$\hat{Y}$} & \textbf{$\hat{Y} - Y$} \\
\hline
2 & 4 & 0 & -4 \\
3 & 6 & 0 & -6 \\
4 & 8 & 0 & -8 \\
\hline
\end{tabular}
\end{center}

\textbf{2. Calcul des dérivées partielles :}

\begin{align*}
\frac{\partial MSE}{\partial a} &= \frac{2}{n}\sum_{i=1}^{n}X_i(\hat{Y}_i - Y_i) \\
                                 &= \frac{2}{3}[2 \times (-4) + 3 \times (-6) + 4 \times (-8)] \\
                                 &= \frac{2}{3}[-8 - 18 - 32] \\
                                 &= \frac{2}{3} \times (-58) \\
                                 &= -\frac{116}{3} \approx -38.67
\end{align*}

\begin{align*}
\frac{\partial MSE}{\partial b} &= \frac{2}{n}\sum_{i=1}^{n}(\hat{Y}_i - Y_i) \\
                                 &= \frac{2}{3}[(-4) + (-6) + (-8)] \\
                                 &= \frac{2}{3} \times (-18) \\
                                 &= -12
\end{align*}

\textbf{3. Mise à jour des paramètres :}

\begin{align*}
a_{new} &= a_{old} - \alpha \frac{\partial MSE}{\partial a} \\
        &= 0 - 0.01 \times (-38.67) \\
        &= 0.3867
\end{align*}

\begin{align*}
b_{new} &= b_{old} - \alpha \frac{\partial MSE}{\partial b} \\
        &= 0 - 0.01 \times (-12) \\
        &= 0.12
\end{align*}

\textbf{4. Nouvelle MSE :}

Nouvelles prédictions avec $a = 0.3867$ et $b = 0.12$ :

\begin{center}
\begin{tabular}{|c|c|c|c|}
\hline
\textbf{X} & \textbf{Y} & \textbf{$\hat{Y}$} & \textbf{$(Y - \hat{Y})^2$} \\
\hline
2 & 4 & 0.8934 & 9.65 \\
3 & 6 & 1.2801 & 22.30 \\
4 & 8 & 1.6668 & 40.11 \\
\hline
\end{tabular}
\end{center}

\begin{align*}
MSE_{new} &= \frac{1}{3}[9.65 + 22.30 + 40.11] \\
          &= \frac{72.06}{3} \\
          &= 24.02
\end{align*}

\textbf{5. Explication de la convergence :}

Le Gradient Descent converge car :
\begin{itemize}
    \item Les dérivées indiquent la direction de la plus forte augmentation de l'erreur
    \item En allant dans le sens opposé (moins la dérivée), on réduit l'erreur
    \item Avec un taux d'apprentissage approprié, on converge vers le minimum
    \item Après plusieurs itérations, $a$ et $b$ se rapprochent des valeurs optimales ($a=2$, $b=0$ dans ce cas)
\end{itemize}
\end{solutionbox}

\begin{tipbox}
\textbf{Tips pour le Gradient Descent :}
\begin{itemize}
    \item Un $\alpha$ trop grand peut faire diverger l'algorithme
    \item Un $\alpha$ trop petit ralentit la convergence
    \item Le Gradient Descent converge vers un minimum local (pas forcément global)
    \item Normalisez vos données pour accélérer la convergence
    \item Utilisez un critère d'arrêt (ex: changement de MSE < seuil)
\end{itemize}
\end{tipbox}

\subsection{Solution Exercice 3 : Questions de cours}

\begin{solutionbox}
\textbf{1. Différence régression simple vs multiple :}
\begin{itemize}
    \item \textbf{Simple :} Une seule variable indépendante : $Y = aX + b$
    \item \textbf{Multiple :} Plusieurs variables indépendantes : $Y = a_1X_1 + a_2X_2 + ... + a_nX_n + b$
\end{itemize}

\textbf{2. MSE comme fonction de coût :}
\begin{itemize}
    \item Pénalise fortement les grandes erreurs (carré)
    \item Toujours positive, facile à minimiser
    \item Différentiable partout
    \item Sensible aux outliers (valeurs aberrantes)
\end{itemize}

\textbf{3. Gradient Descent :}
\begin{itemize}
    \item Algorithme d'optimisation itératif
    \item Calcule le gradient (dérivées) de la fonction de coût
    \item Met à jour les paramètres dans la direction opposée au gradient
    \item Répète jusqu'à convergence
\end{itemize}

\textbf{4. Types de Gradient Descent :}
\begin{itemize}
    \item \textbf{Batch :} Utilise tout le dataset à chaque itération (lent mais précis)
    \item \textbf{Stochastique :} Un échantillon à la fois (rapide mais bruyant)
    \item \textbf{Mini-batch :} Compromis entre les deux (recommandé)
\end{itemize}

\textbf{5. Coefficient $R^2$ :}
$$R^2 = 1 - \frac{SS_{res}}{SS_{tot}} = 1 - \frac{\sum(Y_i - \hat{Y}_i)^2}{\sum(Y_i - \bar{Y})^2}$$
\begin{itemize}
    \item $R^2 = 1$ : Modèle parfait
    \item $R^2 = 0$ : Modèle aussi bon que la moyenne
    \item $R^2 < 0$ : Modèle pire que la moyenne
    \item Mesure la proportion de variance expliquée
\end{itemize}

\textbf{6. Overfitting :}
\begin{itemize}
    \item Le modèle apprend par cœur les données d'entraînement
    \item Faible erreur d'entraînement mais forte erreur de test
    \item Solutions : régularisation (Ridge, Lasso), validation croisée, plus de données
\end{itemize}
\end{solutionbox}

\newpage

% ============================================
% PARTIE 2 : KNN
% ============================================

\section{K-Nearest Neighbors (KNN)}

\subsection{Solution Exercice 1 : Calcul manuel}

\begin{solutionbox}
\textbf{Point de test :} $T = (3, 2)$

\textbf{1. Calcul des distances euclidiennes :}

\begin{align*}
d(T, A) &= \sqrt{(3-1)^2 + (2-2)^2} = \sqrt{4 + 0} = 2.00 \\
d(T, B) &= \sqrt{(3-2)^2 + (2-3)^2} = \sqrt{1 + 1} = 1.41 \\
d(T, C) &= \sqrt{(3-3)^2 + (2-1)^2} = \sqrt{0 + 1} = 1.00 \\
d(T, D) &= \sqrt{(3-4)^2 + (2-2)^2} = \sqrt{1 + 0} = 1.00 \\
d(T, E) &= \sqrt{(3-5)^2 + (2-3)^2} = \sqrt{4 + 1} = 2.24
\end{align*}

\textbf{2. Pour $k=1$ :}

Le plus proche voisin est C ou D (distance = 1.00). Si on choisit C : classe = \textbf{Bleu}

\textbf{3. Pour $k=3$ :}

Les 3 plus proches voisins sont (par ordre croissant de distance) :
\begin{enumerate}
    \item C (distance = 1.00, classe = Bleu)
    \item D (distance = 1.00, classe = Bleu)
    \item B (distance = 1.41, classe = Rouge)
\end{enumerate}

Vote majoritaire : 2 Bleu, 1 Rouge → Classe prédite = \textbf{Bleu}

\textbf{4. Pour $k=5$ :}

Tous les points sont considérés :
\begin{itemize}
    \item Rouge : A, B (2 votes)
    \item Bleu : C, D, E (3 votes)
\end{itemize}

Vote majoritaire : Classe prédite = \textbf{Bleu}

\textbf{5. Influence du choix de $k$ :}
\begin{itemize}
    \item \textbf{$k$ petit (1-3)} : Modèle plus sensible au bruit, frontières complexes
    \item \textbf{$k$ grand} : Modèle plus lisse, frontières simples, peut ignorer les détails locaux
    \item \textbf{$k$ optimal} : Généralement $\sqrt{n}$ où $n$ est le nombre d'échantillons
    \item \textbf{$k$ impair} : Évite les égalités dans le vote majoritaire
\end{itemize}
\end{solutionbox}

\begin{tipbox}
\textbf{Tips pour KNN :}
\begin{itemize}
    \item Normalisez toujours vos données (KNN est sensible aux échelles)
    \item Utilisez $k$ impair pour éviter les égalités
    \item Testez plusieurs valeurs de $k$ avec validation croisée
    \item KNN est coûteux en calcul pour de grands datasets
    \item Considérez des structures de données comme KD-tree pour accélérer
\end{itemize}
\end{tipbox}

\subsection{Solution Exercice 2 : Matrice de confusion}

\begin{solutionbox}
\textbf{1. Matrice de confusion :}

\begin{center}
\begin{tabular}{|c|c|c|}
\hline
 & \textbf{Prédit Positif} & \textbf{Prédit Négatif} \\
\hline
\textbf{Réel Positif} & TP = 3 & FN = 2 \\
\hline
\textbf{Réel Négatif} & FP = 1 & TN = 4 \\
\hline
\end{tabular}
\end{center}

\textbf{Explication des valeurs :}
\begin{itemize}
    \item \textbf{TP = 3} : Points 1, 2, 8 (correctement prédits comme Positifs)
    \item \textbf{TN = 4} : Points 5, 6, 7, 10 (correctement prédits comme Négatifs)
    \item \textbf{FP = 1} : Point 4 (prédit Positif mais réellement Négatif)
    \item \textbf{FN = 2} : Points 3, 9 (prédits Négatifs mais réellement Positifs)
\end{itemize}

\textbf{2. Signification des métriques :}
\begin{itemize}
    \item \textbf{TP (True Positive)} : Vrais positifs - Correctement identifiés comme positifs
    \item \textbf{TN (True Negative)} : Vrais négatifs - Correctement identifiés comme négatifs
    \item \textbf{FP (False Positive)} : Faux positifs - Erreur de type I (alarme fausse)
    \item \textbf{FN (False Negative)} : Faux négatifs - Erreur de type II (manqué)
\end{itemize}

\textbf{3. Calcul des métriques :}

\begin{align*}
\text{Accuracy} &= \frac{TP + TN}{TP + TN + FP + FN} \\
               &= \frac{3 + 4}{3 + 4 + 1 + 2} \\
               &= \frac{7}{10} = 0.70 = 70\%
\end{align*}

\begin{align*}
\text{Precision} &= \frac{TP}{TP + FP} \\
                 &= \frac{3}{3 + 1} \\
                 &= \frac{3}{4} = 0.75 = 75\%
\end{align*}

\begin{align*}
\text{Recall (Sensitivity)} &= \frac{TP}{TP + FN} \\
                            &= \frac{3}{3 + 2} \\
                            &= \frac{3}{5} = 0.60 = 60\%
\end{align*}

\begin{align*}
\text{Specificity} &= \frac{TN}{TN + FP} \\
                  &= \frac{4}{4 + 1} \\
                  &= \frac{4}{5} = 0.80 = 80\%
\end{align*}

\begin{align*}
\text{F1-Score} &= \frac{2 \times Precision \times Recall}{Precision + Recall} \\
               &= \frac{2 \times 0.75 \times 0.60}{0.75 + 0.60} \\
               &= \frac{0.90}{1.35} \\
               &= 0.667 = 66.7\%
\end{align*}

\textbf{4. Interprétation dans un contexte médical :}
\begin{itemize}
    \item \textbf{Precision (75\%)} : Sur 4 diagnostics positifs, 3 sont corrects (1 faux positif)
    \item \textbf{Recall (60\%)} : Sur 5 malades réels, on en détecte 3 (2 manqués - grave !)
    \item \textbf{Specificity (80\%)} : Sur 5 personnes saines, on en identifie correctement 4
    \item \textbf{Accuracy (70\%)} : 70\% des diagnostics sont corrects
\end{itemize}

\textbf{5. Quand privilégier Precision vs Recall :}
\begin{itemize}
    \item \textbf{Privilégier Precision} : Quand les faux positifs sont coûteux
    \begin{itemize}
        \item Exemple : Spam detection (mieux vaut laisser passer un spam que bloquer un email important)
        \item Exemple : Recommandation de produits (mieux vaut ne pas recommander que recommander un mauvais produit)
    \end{itemize}
    \item \textbf{Privilégier Recall} : Quand les faux négatifs sont dangereux
    \begin{itemize}
        \item Exemple : Détection de cancer (mieux vaut un faux positif qu'un faux négatif)
        \item Exemple : Détection de fraude bancaire
    \end{itemize}
\end{itemize}
\end{solutionbox}

\begin{importantbox}
\textbf{Important - Matrice de confusion :}
\begin{itemize}
    \item TP + FN = Total des vrais positifs dans les données
    \item TN + FP = Total des vrais négatifs dans les données
    \item TP + FP = Total des prédictions positives
    \item TN + FN = Total des prédictions négatives
    \item TP + TN + FP + FN = Nombre total d'échantillons
\end{itemize}
\end{importantbox}

\begin{tipbox}
\textbf{Tips pour les métriques :}
\begin{itemize}
    \item Accuracy peut être trompeuse avec des classes déséquilibrées
    \item F1-Score est la moyenne harmonique de Precision et Recall
    \item Utilisez la courbe ROC pour visualiser le compromis Precision/Recall
    \item En cas de classes déséquilibrées, utilisez F1-Score plutôt qu'Accuracy
\end{itemize}
\end{tipbox}

\subsection{Solution Exercice 3 : Questions de cours}

\begin{solutionbox}
\textbf{1. Principe de KNN :}
\begin{itemize}
    \item Algorithme de classification/regression non-paramétrique
    \item Pour classifier un point, trouve les $k$ plus proches voisins
    \item Vote majoritaire (classification) ou moyenne (régression)
    \item Pas d'étape d'entraînement (lazy learning)
\end{itemize}

\textbf{2. Avantages et inconvénients :}
\begin{itemize}
    \item \textbf{Avantages :}
    \begin{itemize}
        \item Simple à comprendre et implémenter
        \item Pas d'hypothèse sur la distribution des données
        \item Efficace pour datasets non-linéaires
        \item Pas d'entraînement nécessaire
    \end{itemize}
    \item \textbf{Inconvénients :}
    \begin{itemize}
        \item Coûteux en calcul (doit calculer toutes les distances)
        \item Sensible aux features non pertinentes
        \item Sensible aux outliers
        \item Nécessite beaucoup de mémoire (stocke tout le dataset)
    \end{itemize}
\end{itemize}

\textbf{3. Choix de $k$ :}
\begin{itemize}
    \item Utiliser validation croisée
    \item Règle empirique : $k = \sqrt{n}$ où $n$ est le nombre d'échantillons
    \item $k$ impair pour éviter les égalités
    \item $k$ trop petit : surapprentissage
    \item $k$ trop grand : sous-apprentissage
\end{itemize}

\textbf{4. Normalisation :}
\begin{itemize}
    \item KNN utilise des distances, donc sensible aux échelles
    \item Normaliser (StandardScaler) pour que toutes les features aient la même importance
    \item Sinon, les features avec grandes valeurs dominent
\end{itemize}

\textbf{5. Distances :}
\begin{itemize}
    \item \textbf{Euclidienne :} $d = \sqrt{\sum(x_i - y_i)^2}$ (sphérique)
    \item \textbf{Manhattan :} $d = \sum|x_i - y_i|$ (en damier)
    \item \textbf{Minkowski :} $d = (\sum|x_i - y_i|^p)^{1/p}$ (généralisation)
    \item Manhattan plus robuste aux outliers
\end{itemize}

\textbf{6. Lazy learner :}
\begin{itemize}
    \item Pas d'étape d'entraînement (pas de modèle appris)
    \item Tout le calcul se fait au moment de la prédiction
    \item Opposé aux "eager learners" (SVM, arbres de décision)
\end{itemize}

\textbf{7. Malédiction de la dimensionnalité :}
\begin{itemize}
    \item En haute dimension, tous les points sont équidistants
    \item La notion de "voisin proche" perd son sens
    \item Solution : réduction de dimensionnalité (PCA, feature selection)
\end{itemize}
\end{solutionbox}

\newpage

% ============================================
% PARTIE 3 : K-MEANS
% ============================================

\section{K-Means Clustering}

\subsection{Solution Exercice 1 : Calcul manuel}

\begin{solutionbox}
\textbf{Points :} P1(1,2), P2(2,3), P3(3,1), P4(8,7), P5(9,8), P6(7,9)

\textbf{Centroïdes initiaux :} C1(2,2), C2(8,8)

\textbf{Itération 1 - Assignment :}

Calcul des distances de chaque point aux centroïdes :

\begin{align*}
d(P1, C1) &= \sqrt{(1-2)^2 + (2-2)^2} = 1.00 \\
d(P1, C2) &= \sqrt{(1-8)^2 + (2-8)^2} = \sqrt{49 + 36} = 9.22 \\
&\Rightarrow P1 \in \text{Cluster 1}
\end{align*}

\begin{align*}
d(P2, C1) &= \sqrt{(2-2)^2 + (3-2)^2} = 1.00 \\
d(P2, C2) &= \sqrt{(2-8)^2 + (3-8)^2} = \sqrt{36 + 25} = 7.81 \\
&\Rightarrow P2 \in \text{Cluster 1}
\end{align*}

\begin{align*}
d(P3, C1) &= \sqrt{(3-2)^2 + (1-2)^2} = \sqrt{1 + 1} = 1.41 \\
d(P3, C2) &= \sqrt{(3-8)^2 + (1-8)^2} = \sqrt{25 + 49} = 8.60 \\
&\Rightarrow P3 \in \text{Cluster 1}
\end{align*}

\begin{align*}
d(P4, C1) &= \sqrt{(8-2)^2 + (7-2)^2} = \sqrt{36 + 25} = 7.81 \\
d(P4, C2) &= \sqrt{(8-8)^2 + (7-8)^2} = 1.00 \\
&\Rightarrow P4 \in \text{Cluster 2}
\end{align*}

\begin{align*}
d(P5, C1) &= \sqrt{(9-2)^2 + (8-2)^2} = \sqrt{49 + 36} = 9.22 \\
d(P5, C2) &= \sqrt{(9-8)^2 + (8-8)^2} = 1.00 \\
&\Rightarrow P5 \in \text{Cluster 2}
\end{align*}

\begin{align*}
d(P6, C1) &= \sqrt{(7-2)^2 + (9-2)^2} = \sqrt{25 + 49} = 8.60 \\
d(P6, C2) &= \sqrt{(7-8)^2 + (9-8)^2} = \sqrt{1 + 1} = 1.41 \\
&\Rightarrow P6 \in \text{Cluster 2}
\end{align*}

\textbf{Itération 1 - Update :}

\begin{align*}
C1_{new} &= \left(\frac{1+2+3}{3}, \frac{2+3+1}{3}\right) = (2, 2) \\
C2_{new} &= \left(\frac{8+9+7}{3}, \frac{7+8+9}{3}\right) = (8, 8)
\end{align*}

Les centroïdes n'ont pas changé ! \textbf{Convergence atteinte.}

\textbf{Clusters finaux :}
\begin{itemize}
    \item \textbf{Cluster 1 :} P1, P2, P3 (centroïde : (2, 2))
    \item \textbf{Cluster 2 :} P4, P5, P6 (centroïde : (8, 8))
\end{itemize}

\textbf{Calcul de l'inertie :}

\begin{align*}
Inertie &= \sum_{i=1}^{2}\sum_{x \in C_i}||x - \mu_i||^2 \\
        &= [d(P1,C1)^2 + d(P2,C1)^2 + d(P3,C1)^2] + [d(P4,C2)^2 + d(P5,C2)^2 + d(P6,C2)^2] \\
        &= [1^2 + 1^2 + 1.41^2] + [1^2 + 1^2 + 1.41^2] \\
        &= [1 + 1 + 2] + [1 + 1 + 2] \\
        &= 4 + 4 = 8
\end{align*}

\textbf{Si $k=3$ :}
\begin{itemize}
    \item On aurait 3 clusters au lieu de 2
    \item L'inertie serait probablement plus petite (meilleur ajustement)
    \item Mais risque de surapprentissage (overfitting)
    \item Utiliser la méthode du coude pour choisir $k$ optimal
\end{itemize}
\end{solutionbox}

\begin{tipbox}
\textbf{Tips pour K-means :}
\begin{itemize}
    \item L'initialisation aléatoire peut donner des résultats différents
    \item Utilisez K-means++ pour une meilleure initialisation
    \item Normalisez vos données avant clustering
    \item K-means suppose des clusters sphériques et de taille similaire
    \item L'inertie diminue toujours quand $k$ augmente (attention au surapprentissage)
\end{itemize}
\end{tipbox}

\subsection{Solution Exercice 2 : Code Python}

\begin{solutionbox}
\textbf{1. Fonction distance euclidienne :}

\begin{lstlisting}
import numpy as np

def euclidean_distance(point1, point2):
    """Calcule la distance euclidienne entre deux points"""
    return np.sqrt(np.sum((point1 - point2) ** 2))
\end{lstlisting}

\textbf{2. Fonction d'assignment :}

\begin{lstlisting}
def assign_clusters(data, centroids):
    """Assigne chaque point au cluster le plus proche"""
    clusters = []
    for point in data:
        distances = [euclidean_distance(point, centroid) 
                     for centroid in centroids]
        cluster = np.argmin(distances)
        clusters.append(cluster)
    return np.array(clusters)
\end{lstlisting}

\textbf{3. Fonction de mise à jour des centroïdes :}

\begin{lstlisting}
def update_centroids(data, clusters, k):
    """Update centroids as mean of points"""
    centroids = []
    for i in range(k):
        cluster_points = data[clusters == i]
        if len(cluster_points) > 0:
            centroid = np.mean(cluster_points, axis=0)
        else:
            # Avoid empty clusters
            centroid = data[np.random.randint(len(data))]
        centroids.append(centroid)
    return np.array(centroids)
\end{lstlisting}

\textbf{4. Algorithme K-means complet :}

\begin{lstlisting}
def kmeans(data, k, max_iter=100, tol=1e-4):
    """Complete K-means implementation"""
    # Random initialization
    n_samples, n_features = data.shape
    centroids = data[np.random.choice(n_samples, k, replace=False)]
    
    for iteration in range(max_iter):
        # Assignment
        clusters = assign_clusters(data, centroids)
        
        # Update
        new_centroids = update_centroids(data, clusters, k)
        
        # Check convergence
        if np.allclose(centroids, new_centroids, atol=tol):
            print(f"Convergence after {iteration+1} iterations")
            break
        
        centroids = new_centroids
    
    return centroids, clusters

# Usage example
data = np.array([[1, 2], [2, 3], [3, 1], [8, 7], [9, 8], [7, 9]])
centroids, clusters = kmeans(data, k=2)
print("Final centroids:", centroids)
print("Clusters:", clusters)
\end{lstlisting}

\textbf{5. Comparaison avec sklearn :}

\begin{lstlisting}
from sklearn.cluster import KMeans

# With sklearn
kmeans_sklearn = KMeans(n_clusters=2, random_state=42, n_init=10)
clusters_sklearn = kmeans_sklearn.fit_predict(data)
centroids_sklearn = kmeans_sklearn.cluster_centers_

print("Sklearn centroids:", centroids_sklearn)
print("Sklearn clusters:", clusters_sklearn)
print("Inertia:", kmeans_sklearn.inertia_)
\end{lstlisting}
\end{solutionbox}

\begin{tipbox}
\textbf{Tips pour le code :}
\begin{itemize}
    \item Utilisez numpy pour des calculs vectorisés (plus rapide)
    \item Gérer les clusters vides (réinitialiser aléatoirement)
    \item Utiliser un critère de tolérance pour la convergence
    \item Fixer random\_state pour la reproductibilité
    \item n\_init dans sklearn teste plusieurs initialisations
\end{itemize}
\end{tipbox}

\subsection{Solution Exercice 3 : Questions de cours}

\begin{solutionbox}
\textbf{1. Principe de K-means :}
\begin{enumerate}
    \item Initialiser $k$ centroïdes aléatoirement
    \item \textbf{Assignment :} Assigner chaque point au cluster le plus proche
    \item \textbf{Update :} Recalculer les centroïdes comme moyenne des points
    \item Répéter 2-3 jusqu'à convergence
\end{enumerate}

\textbf{2. Avantages et inconvénients :}
\begin{itemize}
    \item \textbf{Avantages :}
    \begin{itemize}
        \item Simple et rapide
        \item Efficace pour grands datasets
        \item Garantit la convergence
    \end{itemize}
    \item \textbf{Inconvénients :}
    \begin{itemize}
        \item Nécessite de connaître $k$ à l'avance
        \item Sensible à l'initialisation
        \item Suppose des clusters sphériques
        \item Sensible aux outliers
    \end{itemize}
\end{itemize}

\textbf{3. Choix de $k$ - Méthode du coude :}
\begin{itemize}
    \item Tracer l'inertie en fonction de $k$
    \item Le "coude" (point d'inflexion) indique le $k$ optimal
    \item Autres méthodes : Silhouette score, Gap statistic
\end{itemize}

\textbf{4. Initialisation :}
\begin{itemize}
    \item Initialisation aléatoire peut donner des résultats différents
    \item K-means++ : choisit les centroïdes initiaux intelligemment
    \item Meilleure initialisation = convergence plus rapide et meilleurs résultats
\end{itemize}

\textbf{5. Inertie :}
\begin{itemize}
    \item Somme des distances au carré entre points et centroïdes
    \item Mesure de compacité des clusters
    \item Plus petite = meilleur clustering
    \item Diminue toujours quand $k$ augmente
\end{itemize}

\textbf{6. Clusters vides :}
\begin{itemize}
    \item Peuvent apparaître si un centroïde est trop éloigné
    \item Solutions : réinitialiser, utiliser K-means++, augmenter $k$
\end{itemize}

\textbf{7. Clustering hiérarchique vs K-means :}
\begin{itemize}
    \item \textbf{K-means :} $k$ fixe, partitionnement plat
    \item \textbf{Hiérarchique :} Structure arborescente, pas besoin de $k$
    \item Hiérarchique plus lent mais plus flexible
\end{itemize}

\textbf{8. Hypothèses de K-means :}
\begin{itemize}
    \item Clusters sphériques (distance euclidienne)
    \item Clusters de taille similaire
    \item Variance similaire dans chaque cluster
    \item Si ces hypothèses ne sont pas respectées, K-means peut échouer
\end{itemize}
\end{solutionbox}

\newpage

% ============================================
% PARTIE 4 : ARBRE DE DÉCISION
% ============================================

\section{Arbre de Décision}

\subsection{Solution Exercice 1 : Entropie et gain d'information}

\begin{solutionbox}
\textbf{1. Entropie initiale :}

Total : 14 jours
\begin{itemize}
    \item Oui : 9 jours
    \item Non : 5 jours
\end{itemize}

\begin{align*}
\text{Entropie}(S) &= -p_{\text{Oui}} \log_2(p_{\text{Oui}}) - p_{\text{Non}} \log_2(p_{\text{Non}}) \\
            &= -\frac{9}{14} \log_2\left(\frac{9}{14}\right) - \frac{5}{14} \log_2\left(\frac{5}{14}\right) \\
            &= -0.643 \times (-0.637) - 0.357 \times (-1.485) \\
            &= 0.410 + 0.530 \\
            &= 0.940
\end{align*}

\textbf{2. Entropie pour chaque attribut :}

\textbf{Attribut Outlook :}
\begin{itemize}
    \item \textbf{Ensoleillé :} 5 jours (2 Oui, 3 Non)
    \begin{align*}
    \text{Entropie}(\text{Ensoleillé}) &= -\frac{2}{5}\log_2\left(\frac{2}{5}\right) - \frac{3}{5}\log_2\left(\frac{3}{5}\right) \\
                         &= 0.971
    \end{align*}
    \item \textbf{Nuageux :} 4 jours (4 Oui, 0 Non)
    \begin{align*}
    \text{Entropie}(\text{Nuageux}) &= -\frac{4}{4}\log_2\left(\frac{4}{4}\right) - \frac{0}{4}\log_2\left(\frac{0}{4}\right) \\
                      &= 0 \text{ (cluster pur)}
    \end{align*}
    \item \textbf{Pluvieux :} 5 jours (3 Oui, 2 Non)
    \begin{align*}
    \text{Entropie}(\text{Pluvieux}) &= -\frac{3}{5}\log_2\left(\frac{3}{5}\right) - \frac{2}{5}\log_2\left(\frac{2}{5}\right) \\
                       &= 0.971
    \end{align*}
\end{itemize}

\begin{align*}
\text{Entropie}(S, \text{Outlook}) &= \frac{5}{14} \times 0.971 + \frac{4}{14} \times 0 + \frac{5}{14} \times 0.971 \\
                      &= 0.347 + 0 + 0.347 \\
                      &= 0.694
\end{align*}

\textbf{Attribut Température :}
\begin{itemize}
    \item \textbf{Chaude :} 4 jours (2 Oui, 2 Non) → Entropie = 1.000
    \item \textbf{Douce :} 6 jours (4 Oui, 2 Non) → Entropie = 0.918
    \item \textbf{Fraîche :} 4 jours (3 Oui, 1 Non) → Entropie = 0.811
\end{itemize}

\begin{align*}
\text{Entropie}(S, \text{Température}) &= \frac{4}{14} \times 1.000 + \frac{6}{14} \times 0.918 + \frac{4}{14} \times 0.811 \\
                          &= 0.286 + 0.393 + 0.232 \\
                          &= 0.911
\end{align*}

\textbf{Attribut Humidité :}
\begin{itemize}
    \item \textbf{Élevée :} 7 jours (3 Oui, 4 Non) → Entropie = 0.985
    \item \textbf{Normale :} 7 jours (6 Oui, 1 Non) → Entropie = 0.592
\end{itemize}

\begin{align*}
\text{Entropie}(S, \text{Humidité}) &= \frac{7}{14} \times 0.985 + \frac{7}{14} \times 0.592 \\
                      &= 0.493 + 0.296 \\
                      &= 0.789
\end{align*}

\textbf{3. Gain d'information :}

\begin{align*}
\text{Gain}(S, \text{Outlook}) &= \text{Entropie}(S) - \text{Entropie}(S, \text{Outlook}) \\
                  &= 0.940 - 0.694 = 0.246
\end{align*}

\begin{align*}
\text{Gain}(S, \text{Température}) &= 0.940 - 0.911 = 0.029
\end{align*}

\begin{align*}
\text{Gain}(S, \text{Humidité}) &= 0.940 - 0.789 = 0.151
\end{align*}

\textbf{4. Choix de la racine :}

\textbf{Outlook} a le gain d'information le plus élevé (0.246) → \textbf{Racine = Outlook}

\textbf{5. Construction de l'arbre :}

\begin{verbatim}
                    Outlook
                   /   |   \
          Ensoleillé Nuageux Pluvieux
             /         |         \
          [Oui]    Température  Humidité
                    /    \        /    \
                 Chaude Douce  Élevée Normale
                  /      |       |       |
                [Non]  [Oui]   [Non]   [Oui]
\end{verbatim}

\textbf{6. Classification du nouveau jour :}

Jour : (Ensoleillé, Chaude, Normale)
\begin{enumerate}
    \item Racine : Outlook = Ensoleillé
    \item Sous-arbre Ensoleillé : Température = Chaude → \textbf{Jouer = Non}
\end{enumerate}
\end{solutionbox}

\begin{tipbox}
\textbf{Tips pour les arbres de décision :}
\begin{itemize}
    \item L'entropie est maximale (1.0) quand les classes sont équilibrées (50/50)
    \item L'entropie est minimale (0.0) quand une classe domine (cluster pur)
    \item Le gain d'information mesure la réduction d'incertitude
    \item Choisir toujours l'attribut avec le gain maximum
    \item Arrêter quand l'entropie = 0 (feuille pure) ou plus d'attributs
\end{itemize}
\end{tipbox}

\subsection{Solution Exercice 2 : Gini Impurity}

\begin{solutionbox}
\textbf{1. Gini initial :}

\begin{align*}
Gini(S) &= 1 - p_{Oui}^2 - p_{Non}^2 \\
        &= 1 - \left(\frac{9}{14}\right)^2 - \left(\frac{5}{14}\right)^2 \\
        &= 1 - 0.413 - 0.128 \\
        &= 0.459
\end{align*}

\textbf{2. Gini Gain pour chaque attribut :}

\textbf{Attribut Outlook :}
\begin{itemize}
    \item Ensoleillé : $Gini = 1 - \left(\frac{2}{5}\right)^2 - \left(\frac{3}{5}\right)^2 = 0.480$
    \item Nuageux : $Gini = 1 - \left(\frac{4}{4}\right)^2 - \left(\frac{0}{4}\right)^2 = 0$ (pur)
    \item Pluvieux : $Gini = 1 - \left(\frac{3}{5}\right)^2 - \left(\frac{2}{5}\right)^2 = 0.480$
\end{itemize}

\begin{align*}
Gini(S, Outlook) &= \frac{5}{14} \times 0.480 + \frac{4}{14} \times 0 + \frac{5}{14} \times 0.480 \\
                  &= 0.171 + 0 + 0.171 \\
                  &= 0.342
\end{align*}

\begin{align*}
GiniGain(S, Outlook) &= 0.459 - 0.342 = 0.117
\end{align*}

\textbf{Attribut Humidité :}
\begin{itemize}
    \item Élevée : $Gini = 1 - \left(\frac{3}{7}\right)^2 - \left(\frac{4}{7}\right)^2 = 0.490$
    \item Normale : $Gini = 1 - \left(\frac{6}{7}\right)^2 - \left(\frac{1}{7}\right)^2 = 0.245$
\end{itemize}

\begin{align*}
Gini(S, Humidité) &= \frac{7}{14} \times 0.490 + \frac{7}{14} \times 0.245 \\
                  &= 0.245 + 0.123 \\
                  &= 0.368
\end{align*}

\begin{align*}
GiniGain(S, Humidité) &= 0.459 - 0.368 = 0.091
\end{align*}

\textbf{3. Comparaison Entropie vs Gini :}
\begin{itemize}
    \item \textbf{Entropie :} Gain(Outlook) = 0.246, Gain(Humidité) = 0.151
    \item \textbf{Gini :} Gain(Outlook) = 0.117, Gain(Humidité) = 0.091
    \item Les deux critères donnent le même ordre de préférence
    \item Gini est plus rapide à calculer (pas de logarithme)
    \item Entropie est plus sensible aux différences
    \item En pratique, les résultats sont souvent similaires
\end{itemize}
\end{solutionbox}

\begin{importantbox}
\textbf{Important - Entropie vs Gini :}
\begin{itemize}
    \item Les deux mesurent l'impureté
    \item Entropie : $-\sum p_i \log_2(p_i)$ (plus coûteux)
    \item Gini : $1 - \sum p_i^2$ (plus rapide)
    \item Les deux donnent généralement des arbres similaires
    \item Gini favorise légèrement les splits équilibrés
\end{itemize}
\end{importantbox}

\subsection{Solution Exercice 3 : Questions de cours}

\begin{solutionbox}
\textbf{1. Principe de construction :}
\begin{enumerate}
    \item Choisir l'attribut avec le gain d'information maximum
    \item Créer un nœud pour cet attribut
    \item Diviser le dataset selon les valeurs de l'attribut
    \item Répéter récursivement pour chaque sous-ensemble
    \item Arrêter quand entropie = 0 ou plus d'attributs
\end{enumerate}

\textbf{2. Entropie vs Gini :}
\begin{itemize}
    \item \textbf{Entropie :} Mesure l'incertitude (information)
    \item \textbf{Gini :} Mesure l'impureté (probabilité d'erreur)
    \item Gini plus rapide (pas de log)
    \item Résultats souvent similaires
\end{itemize}

\textbf{3. Gain d'information :}
\begin{itemize}
    \item Mesure la réduction d'incertitude
    \item Gain = Entropie(avant) - Entropie(après)
    \item Plus le gain est élevé, meilleur est le split
    \item Permet de choisir le meilleur attribut à chaque étape
\end{itemize}

\textbf{4. Overfitting :}
\begin{itemize}
    \item L'arbre devient trop profond
    \item Apprend par cœur les données d'entraînement
    \item Mauvaise généralisation
    \item Solutions : pruning, limitation de profondeur, min\_samples\_split
\end{itemize}

\textbf{5. Techniques de pruning :}
\begin{itemize}
    \item \textbf{Pre-pruning :} Arrêter avant (max\_depth, min\_samples)
    \item \textbf{Post-pruning :} Construire l'arbre complet puis élaguer
    \item Post-pruning généralement meilleur mais plus coûteux
\end{itemize}

\textbf{6. Pré vs Post-élagage :}
\begin{itemize}
    \item \textbf{Pre-pruning :} Plus rapide, peut sous-apprendre
    \item \textbf{Post-pruning :} Plus précis, peut mieux généraliser
    \item Post-pruning permet de voir l'arbre complet avant décision
\end{itemize}

\textbf{7. Valeurs manquantes :}
\begin{itemize}
    \item Utiliser la valeur la plus fréquente
    \item Utiliser la valeur moyenne (régression)
    \item Créer une branche "manquante"
    \item Utiliser des algorithmes comme C4.5 qui gèrent les valeurs manquantes
\end{itemize}

\textbf{8. Avantages et inconvénients :}
\begin{itemize}
    \item \textbf{Avantages :}
    \begin{itemize}
        \item Facile à interpréter
        \item Pas besoin de normalisation
        \item Gère les features non-linéaires
        \item Gère les valeurs manquantes
    \end{itemize}
    \item \textbf{Inconvénients :}
    \begin{itemize}
        \item Sujet au surapprentissage
        \item Instable (petit changement → grand changement)
        \item Biais vers les features avec beaucoup de valeurs
    \end{itemize}
\end{itemize}

\textbf{9. Random Forest :}
\begin{itemize}
    \item Ensemble de plusieurs arbres de décision
    \item Chaque arbre entraîné sur un sous-échantillon (bootstrap)
    \item Vote majoritaire pour la prédiction finale
    \item Réduit le surapprentissage et améliore la robustesse
    \item Plus précis mais moins interprétable qu'un seul arbre
\end{itemize}
\end{solutionbox}

\begin{tipbox}
\textbf{Tips pour les arbres de décision :}
\begin{itemize}
    \item Utilisez max\_depth pour limiter la profondeur
    \item min\_samples\_split pour éviter les splits sur peu de données
    \item min\_samples\_leaf pour garantir un minimum dans chaque feuille
    \item Utilisez la validation croisée pour choisir les hyperparamètres
    \item Visualisez l'arbre pour comprendre les décisions
    \item Random Forest est souvent meilleur qu'un seul arbre
\end{itemize}
\end{tipbox}

\newpage

% ============================================
% CONCLUSION
% ============================================

\section{Conclusion}

Ces solutions couvrent les concepts fondamentaux du Machine Learning avec des explications détaillées :

\begin{itemize}
    \item \textbf{Régression Linéaire} : Calculs manuels, MSE, Gradient Descent avec itérations
    \item \textbf{KNN} : Distances, classification, matrice de confusion complète avec toutes les métriques
    \item \textbf{K-means} : Algorithme itératif, calculs de distances et centroïdes, code Python
    \item \textbf{Arbre de Décision} : Entropie, gain d'information, construction complète de l'arbre
\end{itemize}

\textbf{Points clés à retenir :}
\begin{enumerate}
    \item Toujours vérifier vos calculs étape par étape
    \item Comprendre l'interprétation des métriques dans le contexte
    \item Pratiquer les calculs manuels pour bien comprendre les algorithmes
    \item Utiliser les tips pour éviter les erreurs courantes
\end{enumerate}

\textbf{Bon courage pour vos révisions et examens !}

\end{document}

